\section{Abstract}

Fraud detection systems face significant challenges including class imbalance, concept drift, and adversarial evasion attacks. This paper presents an agentic fraud detection system built on the MAPE-K (Monitor-Analyze-Plan-Execute using Knowledge) architecture that autonomously adapts to evolving fraud patterns. The system integrates four key components: social media-based data collection and LLM-driven annotation for generating labeled fraud datasets, CTGAN-based data augmentation to address class imbalance by synthesizing minority class samples, adversarial training to enhance model robustness against evasion attacks, and a closed-loop MAPE-K agent framework that monitors drift, plans retraining, and executes adaptive responses. Our system achieves an F1 score of 0.9911, ROC-AUC of 0.9966, and robustness of 98.34\% against adversarial perturbations. The MAPE-K agents successfully detect concept drift via PSI and KS statistics, trigger automatic retraining, and demonstrate measurable learning adaptation under changing data distributions. This work contributes a complete pipeline from data collection to autonomous model adaptation, with particular emphasis on the agentic intelligence layer as the key differentiator from traditional fraud detection approaches.

\keywords{Fraud Detection \and MAPE-K Architecture \and Adversarial Training \and Concept Drift \and CTGAN \and Autonomous Systems}